给定向量数据库 $\mathcal X\subset\mathbb R^d$ 和查询 $q$，搜索任务是在平方欧氏距离下返回 $k$ 个近似最近邻。本文的研究问题是：如何在单张 GPU 的显存约束内构建并搜索 JHQ 压缩索引，以及如何在给定召回率目标下提高查询吞吐。

搜索先选择相关的 IVF 列表，使用主码粗筛其中的向量，再用残差编码对较小的候选集重新评分。精排预算控制进入最后阶段的候选数量。
真值（ground truth，GT）给出原始向量空间中的精确近邻；召回率是其中被找回的比例。

本文保留 JHQ 不改变维度的正交变换和层级编码（\cref{sec:2-3}），研究 GPU 构建、数据布局、搜索执行和精排预算选择。校准不使用 GT，而是比较较小预算与较大预算参考返回的结果，再为固定索引和假设稳定的查询分布复用一个选定倍率。最终召回率通过 GT 评测核查。

评测使用一张 RTX 5090，执行平方 L2 距离的 top-10 搜索。主要跨系统比较使用四个数据集，batch 为 1,000；查找表、内存布局和预算实验覆盖六个数据集。一项小批量研究增加 batch 8/32（\cref{sec:7-smallbatch}）。索引和核心搜索计算常驻 GPU；计时包含搜索调用内部的查询准备与传输。实现范围不含动态更新、主机分页和多 GPU 执行；十亿规模搜索尚未评测。
